\section{Versuchsaufbau und Durchführung}
\label{sec:Durchführung}

\subsection{Invertierender Verstärker}
\label{subsec:InvertierenderVerstaerker}
Zur Untersuchung des Frequenzverhaltens eines invertierenden Verstärkers wird die Schaltung gemäß \autoref{fig:invertierenderVerstaerker} 
auf einem Steckbrett realisiert. Als Operationsverstärker wird ein 8pin LM741 Chip verwendet.
Zunächst werden die Widerstände auf $R_1 = \qty{1}{\kilo\ohm}$ und $R_2 = \qty{100}{\kilo\ohm}$ gesetzt, was einer erwarteten Verstärkung von 
$-R_2/R_1 = -100$ entspricht. 
Mit einem Funktionsgenerator wird eine sinusförmige Eingangsspannung mit konstanter Amplitude 
$U_\text{in} = \qty{300}{\milli\volt}$ über einen Frequenzbereich von $\qty{1}{\hertz}$ bis $\qty{100}{\kilo\hertz}$ eingespeist.
Für jede Frequenz wurde die Ausgangsspannung $U_\text{out}$ gemessen.
Zur Untersuchung des Verstärkungsverhaltens bei unterschiedlichen Parametern wird die Messung anschließend mit den Parametern 
$(R_1 = \qty{1}{\kilo\ohm}, R_2 = \qty{150}{\kilo\ohm}, U_\text{in} = \qty{100}{\milli\volt})$ und 
$(R_1 = \qty{1}{\kilo\ohm}, R_2 = \qty{10}{\kilo\ohm}, U_\text{in} = \qty{100}{\milli\volt})$ wiederholt.
Die Messdaten werden für jede Konfiguration in doppelt-logarithmischen Bode-Diagrammen dargestellt. Dabei wurde der Frequenzverlauf der 
Verstärkung analysiert und mithilfe eines Fits eine charakteristische Verstärkung, die Grenzfrequenz und die Bandbreite bestimmt.

\subsection{Integrator}
\label{subsec:Integrator}
Das Verhalten eines Integrators, gemäß der Schaltung in \autoref{fig:integrator},
wird - analog zum invertierenden Verstärker - mithilfe eines Steckbretts aufgebaut. Der Widerstand wird auf $R = \qty{10}{\kilo\ohm}$
und die Kapazität auf $C = \qty{22}{\nano\farad}$ festgelegt. Eine sinusförmige Eingangsspannung mit konstanter Amplitude
$U_\text{in} = \qty{1}{\volt}$ wird in den Schaltkreis eingespeist, und die Wertepaare $(U_\text{out}, f)$ werden aufgenommen.
Anschließend wird der Schaltkreis qualitativ auf seine Funktionsweise geprüft, indem ein Dreieck- sowie ein Rechtecksignal angelegt werden.
Die Displayanzeige des Oszilloskops wird jeweils als PNG-Datei gespeichert, und die Abbildungen in der Diskussion~\ref{sec:Diskussion}
auf das erwartete Verhalten hin analysiert. Die verwendete Frequenz und Eingangsspannung werden dabei so angepasst,
dass ein möglichst klares Bild auf dem Oszilloskop zu erkennen ist. Die genauen Parameter sind jeweils in den Abbildungen ersichtlich.

\subsection{Differenzierer}
\label{subsec:Differenzierer}
Das Verhalten einer Differenzierers wird genau Analog zum \autoref{subsec:Integrator} des Integrators analysiert.
Es werden jedoch die Parameter $R = \qty{100}{\kilo\ohm}$, $ C = \qty{22}{\nano\farad}$ und $V_{\text{in}} = \qty{90}{\milli\volt}$ verwendet.

\subsection{Schmitt Trigger}
\label{subsec:SchmittTrigger}
Die Schaltskizze in \autoref{fig:schmittTrigger} eines Schmitt-Triggers wird mithilfe von zwei 8-Pin-LM741-Chips
auf dem Steckbrett realisiert. Die verwendeten Parameter lauten $R_1 = \qty{10}{\kilo\ohm}$ und $R_2 = \qty{100}{\kilo\ohm}$.
Außerdem wird eine sinusförmige Signalspannung mit der Frequenz $\qty{1}{\kilo\hertz}$ verwendet.
Die Eingangsspannung wird zunächst auf $\qty{0}{\volt}$ gesetzt und anschließend schrittweise erhöht,
bis das Ausgangssignal gerade auf den High-Pegel umspringt.
Die zugehörigen Werte für $V_\text{in}$ und $V_\text{out}$ werden notiert.

\subsection{Generator}
\label{subsec:Generator}
Ein Generator gemäß \autoref{fig:generator} wird mit den Parametern 
$R_1 = \qty{10}{\kilo\ohm}$, $R_2 = \qty{100}{\kilo\ohm}$, $R_3 = \qty{1}{\kilo\ohm}$ und $C = \qty{1}{\micro\farad}$ betrieben. 
Das System sollte beim Einschalten der Spannungsversorgung für die 8-Pin-LM741-Chips anfangen zu oszillieren,
dies kann jedoch durch mehrfaches Ein- und Ausschalten der Stromversorgung unterstützt werden. 
Das Bild, das auf dem Display des Oszilloskops sichtbar ist und die Spannungen $U_o$ und $U_1$ aus \ref{fig:generator} zeigt, 
wird als PNG-Datei gespeichert.

\subsection{Generator mit variierenden Amplituden}
\label{subsec:GeneratorMitVariierendenAmplituden}
Mit den Widerständen $R = \qty{10}{\kilo\ohm}$, $R_1 = \qty{1}{\kilo\ohm}$, einem Widerstand $10R$, einem Potenziometer mit bis zu $\qty{1}{\kilo\ohm}$, 
einem Kondensator mit einer Kapazität von $C = \qty{100}{\nano\farad}$ und einem Funktionsgenerator, welcher hier eine rechteckige Spannung produziert, 
wird der Aufbau aus \autoref{fig:generatorMitVariierenderAmplitude} reproduziert.
Die Messdaten des Oszilloskops werden als CSV-Datei gespeichert und zusätzlich wird eine PNG-Datei des Oszilloskop Bildes gesichert.
%Die Messdaten des Oszilloskops werden als CSV-Datei gespeichert und zusätzlich wird jeweils eine PNG-Datei bei maximaler und minimaler 
%Potenziometereinstellung gesichert.
